%=====================================================================
\chapter{Learning from Examples}\label{ch:learning}
%=====================================================================
\locator{6}{Part VI --- Learning, Language and Responsible AI}

\begin{outcomes}
By the end of this chapter you will be able to:
\begin{tight}
  \item \textbf{Name} the three learning paradigms and \textbf{give} a delivery
        example of each.
  \item \textbf{Train} a perceptron by hand on a small table, applying the learning
        rule step by step.
  \item \textbf{Show} that a single perceptron cannot represent exclusive-or, and
        \textbf{explain} what that result did to the field.
  \item \textbf{Recognise} non-convergence from a training trace, and
        \textbf{distinguish} it from slow convergence.
  \item \textbf{State} precisely where this course stops and \emph{Machine
        Learning} begins.
\end{tight}
\end{outcomes}

\begin{prereq}
\chref{introduction} --- and specifically its claim that learning is one route to
intelligent behaviour rather than the definition of it. \chref{uncertainty} for the
probabilistic view of evidence. Nothing in Parts~II to~V depends on this chapter.
\end{prereq}

\begin{keyterms}
supervised learning $\bullet$ unsupervised learning $\bullet$ reinforcement learning
$\bullet$ training example $\bullet$ feature $\bullet$ label $\bullet$ hypothesis
$\bullet$ generalisation $\bullet$ perceptron $\bullet$ weight $\bullet$ bias
$\bullet$ learning rate $\bullet$ perceptron learning rule $\bullet$ epoch $\bullet$
decision boundary $\bullet$ linear separability $\bullet$ exclusive-or $\bullet$
artificial neural network
\end{keyterms}

\section{One Chapter, and Why}

Everything so far has been \emph{told} the rules. Constraints, operators, logical
sentences, conditional probability tables --- somebody wrote them down, and the
agent's job was to reason with them.

This chapter is about the other case: the rules are not available, and only examples
are. It exists for three reasons, and it is important to be clear which.

\textbf{The HEC outlines require it.} The 2017 outline names supervised,
unsupervised and reinforcement learning; the 2023 outline names artificial neural
networks and learning from examples. This chapter covers all of them, in named
sections, and the Course Specification traces each clause here.

\textbf{You cannot judge symbolic methods without it.} \chref{expert} argued that
the expert-system industry collapsed on the knowledge acquisition bottleneck, and
that machine learning dissolves it. That argument is empty unless you have seen a
program acquire a rule from data.

\textbf{And it is a signpost.} \emph{Machine Learning} (Semester 6) has this course
as its published prerequisite and owns the subject entirely. This chapter goes as far
as the perceptron and the exclusive-or limit --- the exact point at which the
historical field stalled for fifteen years --- and then stops. Appendix~F states the
boundary chapter by chapter.

\section{Three Kinds of Learning}

\dsfig{%
\begin{tikzpicture}[font=\scriptsize,
  card/.style={rounded corners=3pt, draw=#1, line width=1.1pt, fill=#1!8,
               text width=4.3cm, align=flush left, inner sep=6pt,
               text=#1!50!black, minimum height=3.9cm},
  hd/.style={rounded corners=2pt, fill=#1, text=white,
             font=\scriptsize\bfseries, inner sep=3pt, text width=4.1cm,
             align=center}]

\node[card=primary] (s) at (0,0)
  {\\[2pt]\textbf{Given:} examples, each with the right answer attached.\\[4pt]
   \emph{Delivery example:} two hundred past tasks, each labelled \emph{late} or
   \emph{on time}.\\[4pt]
   \textbf{Wants:} a rule that labels a task it has never seen.\\[4pt]
   \textbf{The catch:} fitting the examples is easy; being right about new ones is
   the whole problem.};

\node[card=goldhl, right=0.5cm of s] (u)
  {\\[2pt]\textbf{Given:} examples with \emph{no} answers.\\[4pt]
   \emph{Delivery example:} six months of tickets, and the question ``what kinds of
   tickets are there?''\\[4pt]
   \textbf{Wants:} structure --- groups, outliers, dimensions.\\[4pt]
   \textbf{The catch:} there is no right answer to check against, so ``did it
   work?'' is itself a research question.};

\node[card=greenhl, right=0.5cm of u] (r)
  {\\[2pt]\textbf{Given:} an environment, and a reward that arrives late.\\[4pt]
   \emph{Delivery example:} an escalation policy, judged by how the quarter went.\\[4pt]
   \textbf{Wants:} a policy --- what to do in each state.\\[4pt]
   \textbf{The catch:} credit assignment. Which of a thousand decisions caused the
   outcome?};

\node[hd=primary] at (s.north) {SUPERVISED};
\node[hd=goldhl]  at (u.north) {UNSUPERVISED};
\node[hd=greenhl] at (r.north) {REINFORCEMENT};

\node[rounded corners=3pt, fill=lightbg, draw=primary!30, line width=0.8pt,
      text=primary!70!black, font=\scriptsize, text width=14.2cm,
      align=flush left, inner sep=6pt] at (4.85,-2.65)
  {\textbf{What all three have in common, and what separates them from Parts II--V.}
   Every method in this book so far produced an answer that was \emph{correct given
   what it was told} --- a plan that works, a proof that holds, a posterior that
   follows. A learned rule comes with no such guarantee. It comes with an
   \emph{estimate} of how often it will be right on data nobody has seen, and
   producing an honest estimate turns out to be harder than producing the rule.
   That is the subject of \emph{Machine Learning}, and it is why that course spends
   two chapters on evaluation before teaching a single algorithm.};
\end{tikzpicture}}%
{The three learning paradigms, with a delivery example of each. This section
discharges the HEC 2017 outline's \emph{Learning} clause; the sections that follow
take the first one further.}{paradigms}

\section{Supervised Learning}

Two hundred tasks have been completed. For each, the team recorded the estimated
effort in days and a clarity score for the requirement, and whether the task
finished late. The question is whether a rule can be extracted that predicts
lateness for the next task.

Here are twelve of them, enough to work by hand:

\smallskip
\noindent\begin{tabular}{@{}L{1.3cm}L{1.8cm}L{1.8cm}L{1.5cm}L{1.3cm}L{1.8cm}L{1.8cm}L{1.5cm}@{}}
\toprule
\textbf{Task} & \textbf{Effort} & \textbf{Clarity} & \textbf{Late?} &
\textbf{Task} & \textbf{Effort} & \textbf{Clarity} & \textbf{Late?} \\
\midrule
1 & 2 & 9 & no & 7 & 9 & 3 & yes \\
2 & 3 & 8 & no & 8 & 8 & 2 & yes \\
3 & 4 & 9 & no & 9 & 10 & 4 & yes \\
4 & 3 & 7 & no & 10 & 9 & 1 & yes \\
5 & 5 & 8 & no & 11 & 7 & 2 & yes \\
6 & 2 & 6 & no & 12 & 8 & 4 & yes \\
\bottomrule
\end{tabular}

\smallskip
Nobody stated a rule. The rule --- if there is one --- has to come out of the table.

\begin{definitionbox}[The perceptron]
A \term{perceptron} computes a weighted sum of its inputs and fires if the sum
exceeds a threshold:
\[
  \hat{y} \;=\;
  \begin{cases}
    1 & \text{if } \sum_{i} w_{i} x_{i} + b > 0 \\
    0 & \text{otherwise}
  \end{cases}
\]
The \term{weights} $w_{i}$ say how much each feature counts and in which direction;
the \term{bias} $b$ shifts the threshold.

Learning means adjusting them. For each example in turn, apply the
\term{perceptron learning rule}:
\[
  w_{i} \;\leftarrow\; w_{i} + \eta\,(y - \hat{y})\, x_{i},
  \qquad
  b \;\leftarrow\; b + \eta\,(y - \hat{y})
\]
where $y$ is the true label and $\eta$ the \term{learning rate}. If the prediction
is right, $y - \hat{y} = 0$ and nothing changes. If it is wrong, every weight moves
in the direction that would have helped, by an amount proportional to the input.

\smallskip
\noindent\emph{After} F.\ Rosenblatt, ``The Perceptron: A Probabilistic Model for
Information Storage and Organization in the Brain'', \emph{Psychological Review}
65, 1958.
\end{definitionbox}

Trained on the twelve tasks with $\eta = 0.1$, starting from all-zero weights, the
perceptron converges after \textbf{three passes through the data and only two weight
updates}, reaching $100\%$ accuracy. The rule it found is
\[
  0.70 \times \mathit{effort} \;-\; 0.60 \times \mathit{clarity} \;>\; 0
\]
which reads: \emph{a task is predicted late when its effort outweighs its clarity},
in roughly a $6{:}7$ ratio. Nobody wrote that down. It was extracted from twelve
rows, and that is the thing \chref{expert}'s knowledge engineer could not do.

\dsfig{%
\begin{tikzpicture}
\begin{axis}[
  width=11.0cm, height=6.0cm,
  xlabel={estimated effort (days)}, ylabel={requirement clarity},
  xlabel style={font=\scriptsize}, ylabel style={font=\scriptsize},
  tick label style={font=\tiny},
  xmin=0, xmax=12, ymin=0, ymax=11,
  grid=both, major grid style={gray!18}, minor grid style={gray!8},
  legend style={font=\tiny, at={(0.97,0.97)}, anchor=north east,
                draw=gray!40, fill=white},
  legend cell align=left]

\addplot[only marks, mark=*, mark size=2.2pt, greenhl] coordinates {
  (2,9) (3,8) (4,9) (3,7) (5,8) (2,6)};
\addlegendentry{finished on time}
\addplot[only marks, mark=square*, mark size=2.2pt, accent] coordinates {
  (9,3) (8,2) (10,4) (9,1) (7,2) (8,4)};
\addlegendentry{finished late}

% boundary: 0.70 e - 0.60 c = 0  ->  c = (0.70/0.60) e
\addplot[primary, line width=1.3pt, domain=0:9.4] {1.1667*x};
\addlegendentry{$0.70\,e - 0.60\,c = 0$, learned}

\node[font=\tiny\itshape, text=primary, anchor=west] at (axis cs:0.4,9.6)
  {learned from the twelve points,};
\node[font=\tiny\itshape, text=primary, anchor=west] at (axis cs:0.4,8.9)
  {in 3 epochs and 2 updates};
\end{axis}
\end{tikzpicture}}%
{The twelve tasks and the boundary the perceptron learned. Everything above the line
is predicted on time, everything below it late.}{perceptron}

\section{Artificial Neural Networks}

A perceptron is a single artificial neuron: inputs, weights, a summation and a
threshold, loosely modelled on a biological cell. An \term{artificial neural
network} connects many of them in layers, so that the outputs of one layer are the
inputs of the next.

The idea is old --- McCulloch and Pitts proposed the artificial neuron in 1943 ---
and the reason it took until the 1980s to become useful is the subject of the next
section. What a layered network buys is the ability to represent decision boundaries
that are not straight lines, and what it costs is a learning rule: the perceptron
rule works because we know the right answer at the output, and in a hidden layer we
do not. Solving that is \term{backpropagation}, and it is Chapter~13 of
\emph{Machine Learning}.

\section{Linear Separability and the Exclusive-Or Limit}

Look again at \dsref{perceptron}. The perceptron found \emph{a straight line}. That
is all a perceptron can ever find, because $\sum w_{i}x_{i} + b = 0$ is the equation
of a line (a plane, a hyperplane). A dataset that some line can separate is
\term{linearly separable}.

Now a small change to the delivery story. Suppose the team observes that tasks go
late when \emph{exactly one} of two things is true --- the effort is large, or the
requirement is unclear --- and that when \emph{both} are true everybody knows it is
hard, plans accordingly, and it lands on time. Four representative tasks:

\smallskip
\noindent\begin{tabular}{@{}L{2.4cm}L{2.4cm}L{2.4cm}L{6.0cm}@{}}
\toprule
\textbf{Effort} & \textbf{Clarity} & \textbf{Late?} & \textbf{Why} \\
\midrule
small (2) & high (9) & no & easy, and understood \\
large (9) & low (2) & no & obviously hard; planned for \\
large (9) & high (9) & yes & looked easy because it was clear \\
small (2) & low (2) & yes & looked easy because it was small \\
\bottomrule
\end{tabular}

\smallskip
That is \term{exclusive-or}, and no straight line separates those four points.

\dsfig{%
\begin{tikzpicture}
\begin{axis}[
  width=10.4cm, height=5.8cm,
  xlabel={estimated effort (days)}, ylabel={requirement clarity},
  xlabel style={font=\scriptsize}, ylabel style={font=\scriptsize},
  tick label style={font=\tiny},
  xmin=0, xmax=12, ymin=0, ymax=12,
  grid=both, major grid style={gray!18}, minor grid style={gray!8},
  legend style={font=\tiny, at={(0.03,0.97)}, anchor=north west,
                draw=gray!40, fill=white},
  legend cell align=left]

\addplot[only marks, mark=*, mark size=3pt, greenhl] coordinates {(2,9) (9,2)};
\addlegendentry{on time}
\addplot[only marks, mark=square*, mark size=3pt, accent] coordinates {(9,9) (2,2)};
\addlegendentry{late}

\addplot[neutral, line width=0.9pt, dash pattern=on 3pt off 2pt,
         domain=0:12] {12-x};
\addplot[neutral, line width=0.9pt, dash pattern=on 3pt off 2pt,
         domain=0:12] {x};
\addplot[neutral, line width=0.9pt, dash pattern=on 3pt off 2pt,
         domain=0:12] {5.5};

\node[font=\tiny\itshape, text=neutral, anchor=west] at (axis cs:0.4,11.2)
  {three candidate lines --- each gets 3 of 4 right};
\node[font=\tiny\bfseries, text=accent, anchor=west] at (axis cs:3.2,5.9)
  {no line does better};
\end{axis}
\end{tikzpicture}}%
{The exclusive-or arrangement. The two \emph{late} tasks sit at opposite corners,
and so do the two \emph{on time} tasks. Any straight line has at least one point on
the wrong side.}{xor}

Run the perceptron on it and it never stops. The lab allows it $100$, then $1{,}000$,
then $10{,}000$ epochs and counts $383$, $3{,}983$ and $39{,}983$ weight updates ---
the count simply grows with the budget, because the weights are still oscillating
when the budget runs out. Nothing in the trace says ``impossible''; it looks exactly
like slow convergence.

\begin{alertbox}[Proving the limit, not merely observing it]
Failing to converge is weak evidence. The learning rate might be wrong, the
initialisation unlucky, the budget too small --- any practitioner has seen all
three.

So the lab settles it by brute force: it sweeps a grid of every line
$w_{0}x + w_{1}y + b = 0$ over a wide range of coefficients and records the best
accuracy any of them achieves. The answer is \textbf{75\%} --- three points out of
four --- and no line in the grid does better.

That is the difference between ``my training run failed'' and ``this hypothesis
class cannot represent this function''. The first is an engineering problem; the
second is a mathematical fact, and no amount of training fixes it. Learning to tell
them apart is worth more than the perceptron rule itself.
\end{alertbox}

\begin{conceptbox}
In 1969 Minsky and Papert published exactly this analysis --- with more rigour and
over a wider class of problems --- in \emph{Perceptrons}. The effect on the field was
severe: funding for neural network research largely stopped, and the line stayed
closed for roughly fifteen years.

Two things are worth knowing about that episode, because it is routinely retold
badly.

\emph{The result is correct.} A single-layer perceptron genuinely cannot represent
exclusive-or, and the lab reproduces the proof.

\emph{The conclusion drawn from it was not.} A network with a hidden layer can
represent exclusive-or easily; what was missing in 1969 was a way to \emph{train}
one, and backpropagation supplied it in the mid-1980s. So the field abandoned a
line of work because a limitation of the simplest case was read as a limitation of
the approach --- which is \chref{introduction}'s recurring pattern, a true result
about reach mistaken for a verdict.
\end{conceptbox}

\section{Unsupervised Learning}

Six months of tickets, no labels, and the question \emph{what kinds of ticket do we
actually get?}

There is no right answer to check against, which is what makes this different in
kind rather than in difficulty. A program can group tickets by similarity and the
groups may be illuminating or useless; nothing in the data says which. Judging the
result requires a person who knows the domain, and inventing measures that stand in
for that judgement is an active research area.

That is as far as this course goes. \emph{Machine Learning} Chapters~15 and~16 give
$k$-means, hierarchical clustering, principal component analysis and
self-organizing maps, with the evaluation machinery that makes them usable.

\section{Reinforcement Learning}

An escalation policy is not judged ticket by ticket. It is judged by how the quarter
went --- and a quarter contains a thousand decisions, of which perhaps three
mattered.

\dsfig{%
\begin{tikzpicture}[font=\scriptsize,
  b/.style={rounded corners=4pt, draw=#1, line width=1.2pt, fill=#1!10,
            text=#1!55!black, font=\small\bfseries, minimum width=3.0cm,
            minimum height=1.0cm, align=center},
  a/.style={-{Stealth[length=2.4mm]}, line width=1.1pt}]

\node[b=primary] (ag) at (0,0)  {agent\\{\scriptsize the escalation policy}};
\node[b=goldhl]  (en) at (6.2,0){environment\\{\scriptsize the queue, the team,
  the quarter}};

\draw[a, accent] (ag.north east) .. controls +(1.6,0.9) and +(-1.6,0.9) ..
  (en.north west) node[midway, above, font=\scriptsize, text=accent]
  {\textbf{action}: escalate, defer, merge};
\draw[a, secondary] (en.south west) .. controls +(-1.6,-0.9) and +(1.6,-0.9) ..
  (ag.south east) node[midway, below, font=\scriptsize, text=secondary]
  {\textbf{state} and \textbf{reward}};

\node[rounded corners=3pt, fill=lightbg, draw=primary!30, line width=0.9pt,
      text=primary!70!black, font=\scriptsize, text width=13.0cm,
      align=flush left, inner sep=6pt] at (3.1,-2.2)
  {\textbf{The hard part is the delay.} The reward for tonight's escalation arrives
   at the end of the quarter, mixed with the consequences of nine hundred and
   ninety-nine other decisions. Working out which actions deserve the credit is the
   \term{credit assignment problem}, and it is the reason reinforcement learning is
   a separate subject rather than a variation on the previous two.

   \smallskip
   Note the shape of this diagram: it is \dsref{agent-loop} from \chref{agents},
   with a reward signal added. The vocabulary of Part~I was chosen to make that
   continuity visible.};
\end{tikzpicture}}%
{The reinforcement learning loop. One figure is all this course gives it;
\emph{Machine Learning} Chapters 18--20 develop Markov decision processes, Monte
Carlo methods and Q-learning.}{rl-loop}

\section{Where This Course Hands Over}

\begin{alertbox}[The boundary, stated exactly]
\textbf{This chapter has taught:} the three paradigms and what distinguishes them; a
perceptron trained by hand; the fact that it finds a straight line; and a proof that
some problems have no straight line.

\textbf{This course does not teach, and \emph{Machine Learning} does:} multi-layer
networks and backpropagation (its Chapter~13); how to tell whether a learned rule
will work on data it has not seen --- generalisation, overfitting, the bias--variance
trade-off, cross-validation (Chapters~7 and~8, taught \emph{before} the algorithms,
deliberately); decision trees, naive Bayes, $k$-nearest neighbours, support vector
machines, ensembles (Chapters~9--14); clustering and dimensionality reduction
(15--17); and Markov decision processes and reinforcement learning (18--20).

\textbf{Why not here.} Two reasons, and neither is that the material is too hard.
It would duplicate a course three semesters later that has the statistical
prerequisites to do it properly --- and the room it would take comes out of the
symbolic material, which no other core course in the degree covers at all. A
graduate who has not met resolution or arc consistency will not meet them anywhere
else; a graduate who has not met backpropagation meets it in Semester~6.
\end{alertbox}

\begin{worked}[Training the perceptron by hand]
Weights start at zero, $\eta = 0.1$, and the rule fires when
$w_{1}e + w_{2}c + b > 0$. Take the examples in order.

\smallskip
\textbf{Task 1} $(e,c) = (2,9)$, not late. Sum $= 0$, which is not $> 0$, so
$\hat{y} = 0$. Correct, and \emph{nothing changes} --- the rule only learns from
mistakes.

\smallskip
Tasks 2--6 are all \emph{not late} and all still predicted $0$. No change.

\smallskip
\textbf{Task 7} $(9,3)$, late. Sum $= 0$, so $\hat{y} = 0$ and
$y - \hat{y} = 1$. Update:
\[
  w_{1} = 0 + 0.1 \times 1 \times 9 = 0.9, \quad
  w_{2} = 0 + 0.1 \times 1 \times 3 = 0.3, \quad
  b = 0 + 0.1 = 0.1
\]
Now the rule is $0.9e + 0.3c + 0.1 > 0$ --- which predicts \emph{everything} late,
since all inputs are positive. It has over-corrected, and that is normal.

\smallskip
\textbf{Tasks 8--12} are all late and now all predicted late. No change. End of the
first pass: \textbf{one error}.

\smallskip
\textbf{Second pass, task 1} $(2,9)$, not late. Sum
$= 0.9(2) + 0.3(9) + 0.1 = 1.8 + 2.7 + 0.1 = 4.6 > 0$, so $\hat{y} = 1$: wrong.
Now $y - \hat{y} = -1$:
\[
  w_{1} = 0.9 - 0.1(2) = 0.7, \quad
  w_{2} = 0.3 - 0.1(9) = -0.6, \quad
  b = 0.1 - 0.1 = 0.0
\]
\textbf{The clarity weight has gone negative}, which is the moment the rule becomes
sensible: more clarity now argues \emph{against} lateness. The learning rule was not
told that; it followed from one mistake on one example.

\smallskip
Tasks 2--12 in the second pass are all now predicted correctly. End of pass two:
\textbf{one error}.

\smallskip
\textbf{Third pass:} zero errors. Converged, after $3$ epochs and $2$ updates, with
\[
  0.70\,e \;-\; 0.60\,c \;>\; 0
\]

\smallskip
\textbf{What to take from it.} Both updates came from a single misclassified
example, and the second one produced the insight --- a negative weight on clarity
--- that a person would have had to think of. That is the appeal of learning, in
twelve rows and two updates.

\smallskip
\textbf{And what to be careful about.} The rule has $100\%$ accuracy \emph{on the
data it was trained on}, which is the least interesting fact about it. Whether it is
right about the next task is a different question, and this course does not have the
machinery to answer it. That question is the first thing \emph{Machine Learning}
takes up.
\end{worked}

\begin{pitfall}
\begin{tight}
  \item \textbf{Reading training accuracy as performance.} $100\%$ on the twelve
        rows it learned from says almost nothing about the thirteenth task.
  \item \textbf{Mistaking non-convergence for slow convergence.} The XOR trace looks
        identical to a run that needs more epochs. Test whether \emph{any}
        hypothesis in the class can fit before tuning anything.
  \item \textbf{Concluding that neural networks cannot do exclusive-or.} A single
        perceptron cannot; a hidden layer can, easily. The 1969 result was correct
        and the inference drawn from it was not.
  \item \textbf{Expecting a learned rule to come with a guarantee.} Parts~II to~V
        produced plans that work and proofs that hold. A learned rule comes with an
        error estimate, and producing an honest one is harder than producing the
        rule.
  \item \textbf{Treating unsupervised learning as supervised learning without
        labels.} There is no right answer to check against, which changes what
        ``it worked'' can even mean.
  \item \textbf{Believing this chapter is a machine learning course.} It is four
        sections and a signpost. Appendix~F lists what is missing.
\end{tight}
\end{pitfall}

%---------------------------------------------------------------------
\begin{lab}[The Perceptron and Its Limit]
The learning rule from scratch on twelve tasks, then on the exclusive-or
arrangement, then the brute-force proof that no line can do it.
\begin{lstlisting}[language=Python]
"""Chapter 17 lab -- the perceptron, and where this course stops.

Two datasets and one proof. The first converges in three passes. The
second never converges, and the grid sweep at the end shows that this
is a fact about straight lines rather than about the training run.
"""

# twelve past tasks: (effort days, requirement clarity) -> late?
TASKS = [
    ((2, 9), 0), ((3, 8), 0), ((4, 9), 0), ((3, 7), 0),
    ((5, 8), 0), ((2, 6), 0),
    ((9, 3), 1), ((8, 2), 1), ((10, 4), 1), ((9, 1), 1),
    ((7, 2), 1), ((8, 4), 1),
]

# the exclusive-or arrangement: late when EXACTLY ONE of
# "large effort" and "low clarity" holds
XOR_TASKS = [
    ((2, 9), 0),   # small effort, high clarity -> on time
    ((9, 2), 0),   # large effort, low clarity  -> on time; they planned
    ((9, 9), 1),   # large effort, high clarity -> late; looked easy
    ((2, 2), 1),   # small effort, low clarity  -> late; looked small
]


def predict(x, w, b):
    return 1 if w[0] * x[0] + w[1] * x[1] + b > 0 else 0


def accuracy(data, w, b):
    return sum(predict(x, w, b) == y for x, y in data) / len(data)


def train(data, eta=0.1, epochs=100, trace=False):
    """The perceptron learning rule. Learns only from mistakes."""
    w, b, updates = [0.0, 0.0], 0.0, 0
    for ep in range(1, epochs + 1):
        errors = 0
        for (x, y) in data:
            err = y - predict(x, w, b)
            if err:
                errors += 1
                updates += 1
                w[0] += eta * err * x[0]
                w[1] += eta * err * x[1]
                b += eta * err
        if trace:
            print("   epoch %2d: %d error(s), w = (%.2f, %.2f), b = %.2f"
                  % (ep, errors, w[0], w[1], b))
        if errors == 0:
            return w, b, ep, updates, True
    return w, b, epochs, updates, False


def best_line_accuracy(data, span=20, step=0.5):
    """Brute force over a grid of ALL lines w0*x + w1*y + b = 0.

    This is the difference between 'training failed' and 'no line
    exists'. It is only possible because the problem is tiny.
    """
    best, arg = 0.0, None
    vals = [i * step for i in range(-span * 2, span * 2 + 1)]
    for w0 in vals:
        for w1 in vals:
            for bb in vals:
                a = accuracy(data, [w0, w1], bb)
                if a > best:
                    best, arg = a, (w0, w1, bb)
    return best, arg


if __name__ == "__main__":
    print("=== twelve past tasks ===")
    w, b, epochs, updates, converged = train(TASKS, trace=True)
    print("   converged: %s, after %d epochs and %d weight updates"
          % (converged, epochs, updates))
    print("   learned rule:  %.2f*effort %+.2f*clarity %+.2f > 0"
          % (w[0], w[1], b))
    print("   accuracy on the training data: %.0f%%"
          % (100 * accuracy(TASKS, w, b)))
    assert converged and epochs == 3 and updates == 2
    assert accuracy(TASKS, w, b) == 1.0
    # the clarity weight went NEGATIVE: more clarity argues against late
    assert w[1] < 0
    print()
    print("   The weight on clarity is negative. Nobody said that;")
    print("   it came from one misclassified example in pass two.")
    print()
    print("   But 100% here is accuracy on the data it LEARNED from.")
    print("   Whether it is right about the next task is a different")
    print("   question, and this course cannot answer it.")

    print()
    print("=== the exclusive-or arrangement ===")
    for (x, y) in XOR_TASKS:
        print("   effort %2d, clarity %2d -> %s"
              % (x[0], x[1], "late" if y else "on time"))

    print()
    print("   epochs allowed   converged   weight updates")
    print("   " + "-" * 44)
    for budget in (100, 1000, 10000):
        _, _, _, ups, conv = train(XOR_TASKS, epochs=budget)
        print("   %14d   %9s   %13d" % (budget, conv, ups))
    _, _, _, _, conv = train(XOR_TASKS, epochs=10000)
    assert not conv
    print()
    print("   The update count just grows with the budget: the")
    print("   weights are still oscillating when time runs out.")
    print("   Nothing here says 'impossible' -- it looks exactly")
    print("   like slow convergence.")

    # so settle it by brute force
    print()
    print("=== can ANY straight line do it? ===")
    best, arg = best_line_accuracy(XOR_TASKS)
    print("   best accuracy over a grid of every line: %.0f%%"
          % (100 * best))
    print("   achieved by %.1f*effort %+.1f*clarity %+.1f > 0"
          % arg)
    assert best == 0.75
    print()
    print("   Three of four, and no line does better. This is not a")
    print("   training failure -- it is a fact about the hypothesis")
    print("   class. No learning rate, initialisation or patience")
    print("   fixes it; only a richer class does, and a hidden layer")
    print("   is exactly that.")

    # and the same sweep on the separable data, for contrast
    best2, _ = best_line_accuracy(TASKS)
    assert best2 == 1.0
    print()
    print("   (The same sweep on the twelve tasks finds a line with")
    print("   100% accuracy, which is why the perceptron found one.)")

    print()
    print("=" * 58)
    print("Machine Learning Fall 2026, Chapter 13, takes over here.")
    print("=" * 58)
\end{lstlisting}
Two things to try. Set \texttt{eta} to $10$ and re-run the first dataset: it still
converges, because the perceptron convergence theorem does not depend on the
learning rate --- only on separability. Then change one label in \texttt{TASKS} to
make the data non-separable, and watch the first dataset start behaving like the
second.
\end{lab}

%---------------------------------------------------------------------
\begin{checkpoint}
\begin{tightnum}
  \item The perceptron made only two weight updates in three passes over twelve
        examples. Explain why so few, and what the second update changed about the
        meaning of the rule.
  \item A colleague's perceptron has run for ten thousand epochs without
        converging. Give two possible explanations and describe the single cheapest
        experiment that distinguishes them.
  \item State what a single perceptron cannot represent, what the 1969 conclusion
        from that result was, and why the conclusion was wrong.
\end{tightnum}
\end{checkpoint}

%---------------------------------------------------------------------
\begin{chaptersummary}
\begin{tight}
  \item \term{Supervised} learning is given labelled examples and wants a rule that
        generalises; \term{unsupervised} learning is given unlabelled data and wants
        structure, with no right answer to check against; \term{reinforcement}
        learning is given delayed reward and must solve \term{credit assignment}.
  \item A \term{perceptron} fires when $\sum w_{i}x_{i} + b > 0$. The
        \term{perceptron learning rule} adjusts weights only on mistakes, by
        $\eta(y - \hat{y})x_{i}$.
  \item On twelve separable tasks it converged in $3$ epochs and $2$ updates to
        $0.70\,e - 0.60\,c > 0$. The negative weight on clarity --- the insight in
        the rule --- came from one misclassified example.
  \item A perceptron can only ever find a \emph{straight line}. The
        \term{exclusive-or} arrangement has no separating line, and the lab proves
        it by sweeping every line and finding that the best reaches $75\%$.
  \item Distinguish \emph{training failed} from \emph{no hypothesis in the class
        fits}. The traces look identical; only the second is a mathematical fact,
        and no tuning fixes it.
  \item Minsky and Papert's 1969 result was correct and the inference drawn from it
        was not: a hidden layer represents exclusive-or easily, and what was missing
        was a way to train one.
  \item This is a signpost, not a unit. Generalisation, evaluation,
        backpropagation, trees, kernels, ensembles, clustering and reinforcement
        learning belong to \emph{Machine Learning} in Semester~6; Appendix~F lists
        the boundary chapter by chapter.
\end{tight}
\end{chaptersummary}

\begin{reviewq}
  \item \textbf{(MCQ)} The perceptron learning rule updates the weights:
        (a)~after every example (b)~only when the prediction is wrong
        (c)~once per epoch (d)~only at the end of training.
  \item \textbf{(MCQ)} A single perceptron can represent exactly the functions that
        are: (a)~monotone (b)~linearly separable (c)~continuous (d)~boolean.
  \item \textbf{(MCQ)} The best accuracy any straight line achieves on the
        exclusive-or arrangement is: (a)~$50\%$ (b)~$75\%$ (c)~$90\%$ (d)~$100\%$.
  \item \textbf{(MCQ)} Credit assignment is the central difficulty of:
        (a)~supervised learning (b)~unsupervised learning (c)~reinforcement
        learning (d)~concept learning.
  \item \textbf{(Short)} Explain why a training run that has not converged is weak
        evidence that the problem is unlearnable, and what would be strong evidence.
  \item \textbf{(Short)} The learned rule scores $100\%$ on its training data. State
        why that is the least interesting fact about it, and name the course that
        addresses the interesting question.
  \item \textbf{(Short)} Give one delivery question that is supervised, one that is
        unsupervised and one that is reinforcement learning, and say what makes each
        one that kind.
  \item \textbf{(Applied)} Work the perceptron by hand on the four exclusive-or
        tasks for two full passes, starting from zero weights with $\eta = 0.1$.
        Show the weights after each update, and identify the point at which the
        sequence begins to repeat.
\end{reviewq}
